\documentclass[10pt,letterpaper,landscape]{article}
\usepackage[spanish]{babel}
\usepackage{fontspec}
\setmainfont{TeX Gyre Heros}
\usepackage[margin=1.25cm]{geometry}
\usepackage{array}
\usepackage{longtable}
\usepackage[table]{xcolor}
\usepackage{hyperref}
\definecolor{azul}{HTML}{173B57}
\definecolor{claro}{HTML}{E9F0F5}
\hypersetup{
  pdftitle={Prototipo - Rúbrica Reporte 01 - Ley de Ohm},
  pdfauthor={Diego García}
}
\setlength{\parindent}{0pt}
\setlength{\parskip}{4pt}
\renewcommand{\arraystretch}{1.25}
\newcolumntype{P}[1]{>{\raggedright\arraybackslash}p{#1}}
\begin{document}

{\color{azul}\LARGE\bfseries Rúbrica prototipo - Reporte 01}\\
{\large Ley de Ohm: voltaje en una resistencia}

\colorbox{red!10}{\parbox{0.98\textwidth}{
\textbf{NO PUBLICAR.} Puntaje diseñado únicamente para probar estructura,
lectura y carga en Canvas. No posee conversión oficial a nota.}}

\small
\begin{longtable}{P{3.0cm}P{1.0cm}P{4.55cm}P{4.35cm}P{4.15cm}P{3.85cm}}
\rowcolor{azul}
\color{white}\textbf{Criterio y evidencia} &
\color{white}\textbf{Máx.} &
\color{white}\textbf{Completo} &
\color{white}\textbf{Mayormente logrado} &
\color{white}\textbf{Parcial} &
\color{white}\textbf{Mínimo o cero}\\
\endfirsthead
\rowcolor{azul}
\color{white}\textbf{Criterio y evidencia} &
\color{white}\textbf{Máx.} &
\color{white}\textbf{Completo} &
\color{white}\textbf{Mayormente logrado} &
\color{white}\textbf{Parcial} &
\color{white}\textbf{Mínimo o cero}\\
\endhead

\rowcolor{claro}
\textbf{1. Registro trazable}\\
Valor nominal, tolerancia, \(R\) medida, rangos y seis pares \(I,V\). &
3 &
\textbf{3,0:} todos los datos, unidades, resolución/rango y observaciones son
trazables. &
\textbf{2,25:} falta un dato secundario o existe una inconsistencia menor. &
\textbf{1,5:} hay datos suficientes para analizar, pero faltan varios metadatos
o unidades. &
\textbf{0,75:} registro fragmentario. \textbf{0:} no hay datos evaluables.\\

\textbf{2. Montaje y uso de instrumentos}\\
Diagrama y explicación de conexiones y seguridad. &
3 &
\textbf{3,0:} voltímetro paralelo, amperímetro serie y ohmímetro desenergizado;
incluye cuidados pertinentes. &
\textbf{2,25:} montaje correcto con una omisión menor en explicación o cuidado. &
\textbf{1,5:} idea general correcta, pero conexión o precaución queda ambigua. &
\textbf{0,75:} errores relevantes sin llegar a un cortocircuito explícito.
\textbf{0:} propone una conexión peligrosa o no presenta montaje.\\

\rowcolor{claro}
\textbf{3. Gráfico \(V\) versus \(I\)}\\
Ejes, unidades, puntos y representación legible. &
4 &
\textbf{4,0:} gráfico completo, escala adecuada, ejes y unidades correctos,
todos los puntos identificables. &
\textbf{3,0:} gráfico interpretable con un detalle menor de rotulación o escala. &
\textbf{2,0:} muestra la tendencia, pero tiene errores de eje, unidad o
presentación. &
\textbf{1,0:} gráfico apenas interpretable. \textbf{0:} ausente o no corresponde
a los datos.\\

\end{longtable}
\newpage
\begin{longtable}{P{3.0cm}P{1.0cm}P{4.55cm}P{4.35cm}P{4.15cm}P{3.85cm}}
\rowcolor{azul}
\color{white}\textbf{Criterio y evidencia} &
\color{white}\textbf{Máx.} &
\color{white}\textbf{Completo} &
\color{white}\textbf{Mayormente logrado} &
\color{white}\textbf{Parcial} &
\color{white}\textbf{Mínimo o cero}\\

\textbf{4. Ajuste e interpretación de la pendiente}\\
Modelo \(V=mI+b\), \(m\), \(b\), unidades y significado físico. &
4 &
\textbf{4,0:} ajuste correcto; pendiente e intercepto con unidades y una medida
de calidad; interpreta \(m\) como resistencia. &
\textbf{3,0:} cálculo e interpretación correctos con una omisión menor. &
\textbf{2,0:} ajuste incompleto o interpretación parcialmente correcta. &
\textbf{1,0:} usa solo \(V/I\) en un punto o comete un error mayor recuperable.
\textbf{0:} no existe análisis evaluable.\\

\rowcolor{claro}
\textbf{5. Comparación y calidad de medición}\\
Compara pendiente, ohmímetro y valor nominal. &
3 &
\textbf{3,0:} comparación cuantitativa, diferencia porcentual y discusión de
tolerancia, resolución, contactos y calentamiento. &
\textbf{2,25:} comparación correcta con discusión limitada. &
\textbf{1,5:} comparación cualitativa o cálculo parcial. &
\textbf{0,75:} enumera valores sin compararlos. \textbf{0:} ausente.\\

\textbf{6. Conclusión y taller contextualizado}\\
Responde la pregunta con evidencia y distingue resistencia metálica de
conductancia electrolítica. &
3 &
\textbf{3,0:} conclusión sustentada en pendiente/intercepto y puente
químico-farmacéutico preciso, con límites del modelo. &
\textbf{2,25:} conclusión correcta y contextualización breve con una omisión. &
\textbf{1,5:} repite la ley de Ohm, usa poca evidencia o generaliza demasiado
la analogía. &
\textbf{0,75:} conclusión vaga. \textbf{0:} ausente o contradice los datos.\\
\end{longtable}

\textbf{Total: 20 puntos de prototipo.}

\textbf{Errores de arrastre:} un error inicial se descuenta en el criterio
donde se origina; los pasos posteriores se evalúan por consistencia cuando el
procedimiento sigue siendo físicamente interpretable.

\textbf{Retroalimentación:} todo descuento debe identificar la evidencia
faltante, el error y una acción concreta de mejora.

\end{document}
